% =============================================================================
% METADATOS DEL DOCUMENTO (plantilla genérica para prácticas de laboratorio)
% Este archivo debe cargarse AL PRINCIPIO del preámbulo.
% =============================================================================
%
% INSTRUCCIONES: Reemplaza los textos entre corchetes por los valores reales
% de tu compendio. Puedes eliminar los comentarios que no necesites.
% =============================================================================

% --- Identificación del documento ---
\newcommand{\docTitulo}{QbitLab: Laboratorio de Física Computacional}
\newcommand{\docSubtitulo}{Prácticas de Machine Learning para Físicos}
\newcommand{\docAutor}{Jesús Antonio Chala Casanova (\texttt{@QbitLab})}
\newcommand{\docFecha}{\the\year}                     % Año actual
\newcommand{\docEstado}{Versión 1.0}                 % Versión del compendio

% --- Cita de portada (opcional) ---
\newcommand{\docCita}{}

% --- Textos para páginas preliminares (opcionales) ---
\newcommand{\docDedicatoria}{}                        % Si no se usa, dejar vacío
\newcommand{\docEpigrafe}{}                           % Si no se usa, dejar vacío
\newcommand{\docAgradecimientos}{}

% --- Palabras clave y clasificación (para metadatos PDF) ---
\newcommand{\docKeywords}{física computacional, machine learning, redes neuronales,
  optimización estocástica, modelos de difusión, ecuaciones diferenciales estocásticas}
\newcommand{\docSubject}{Física Computacional / Machine Learning}

% =============================================================================
% NUEVOS COMANDOS (opcionales, pueden usarse en la portada o en secciones)
% =============================================================================

% --- Afiliación del autor ---
\newcommand{\docAfiliacion}{Departamento de Física, Universidad [Nombre]}

% --- Fechas específicas (si se desea distinguir realización y entrega) ---
\newcommand{\docFechaRealizacion}{\today}
\newcommand{\docFechaEntrega}{\today}

% --- Resumen general del compendio ---
\newcommand{\docAbstract}{%
  Este compendio reúne prácticas de laboratorio diseñadas para introducir a
  estudiantes de física en el machine learning desde una perspectiva computacional.
  Los temas incluyen aproximación de funciones, regularización, optimización
  estocástica, arquitecturas neuronales, modelos generativos y difusión,
  con énfasis en la conexión con la mecánica estadística y la simulación numérica.
  Cada práctica incluye ejercicios computacionales independientes del lenguaje
  de programación, fomentando la implementación en el entorno preferido del
  estudiante.%
}

% --- Referencia a material complementario (opcional) ---
\newcommand{\docLibroReferencia}{}                     % No utilizado en este proyecto

% =============================================================================
% RECURSOS BIBLIOGRÁFICOS
% =============================================================================
\addbibresource{references/references.bib}
\addbibresource{references/ReferencesChapter1.bib}
